\chapter{Construction~: un moteur de risque}\label{ch:rc:build-a-risk-engine}

À six heures du matin, le rapport de risque doit être sur les écrans des responsables de desk~: chaque transaction réévaluée sous
chaque scénario, le P\&L sommé desk par desk, la VaR et l'\omterm{def:rc:market-risk-measures:es}{expected shortfall} à chaque nœud, chaque dépassement de limite signalé. Entre
les données de marché de fin de journée et cette échéance se trouve un traitement par lots qui doit tenir dans sa fenêtre chaque nuit~;
dans une enquête réalisée en 2011 par McKinsey auprès de banques, les durées moyennes de calcul de la VaR allaient de 2 à 15 heures,
et davantage en situation de stress. Ce dernier chapitre construit le \omterm{def:rc:build-a-risk-engine:engine}{moteur de risque} de la firme au-dessus
de la bibliothèque de pricing du livre~5 et des mesures du \cref{ch:rc:market-risk-measures}, et pose la question que tout moteur de ce type
doit affronter~: quelle exactitude l'échéance peut-elle se permettre~?

\section{Architecture~: données de marché, scénarios, réévaluation, agrégation}

\begin{definition}[Moteur de risque, agrégation des données de risque]\label{def:rc:build-a-risk-engine:engine}
Un \emph{moteur de risque}\index{moteur de risque} est le système qui transforme les positions de la firme et les données de marché
en ses mesures de risque~: il construit des scénarios, réévalue les positions sous ces scénarios, agrège les résultats
selon la structure de la firme, et les rapporte par rapport aux limites. L'\emph{agrégation des données de risque}\index{agrégation des données de risque}
est, selon les termes du Comité de Bâle, la définition, la collecte et le traitement des données de risque conformément aux
exigences de reporting de la banque, afin de lui permettre de mesurer sa performance par rapport à sa tolérance ou à son appétit pour le risque.
\end{definition}

Le moteur comporte quatre étapes (\cref{fig:rc:build-a-risk-engine:arch}). Un instantané des données de marché fixe les courbes,
les spots et les volatilités du jour. Les générateurs de scénarios transforment l'historique (ou les scénarios de stress, chapitre~22) en
chocs sur des \omterm{def:rc:market-risk-measures:factor}{facteurs de risque} nommés. La réévaluation transforme chaque position et chaque scénario en P\&L, grâce à la bibliothèque de pricing,
de sorte que le \omterm{def:rc:build-a-risk-engine:engine}{moteur de risque} ne price jamais rien lui-même. L'agrégation somme le P\&L le long de la \omterm{def:rc:build-a-risk-engine:hierarchy}{hiérarchie de risque}
et calcule les mesures. Le tableau des P\&L, scénarios par transactions, est l'objet central du moteur~:
chaque rapport en est une fonction.

\begin{omfigure}
\begin{tikzpicture}[font=\small,box/.style={draw,rounded corners,minimum height=1cm,align=center,text width=2.3cm}]
\node[box,fill=omThm!10] (md) at (0,1.4) {instantané de données de marché};
\node[box,fill=omThm!10] (tr) at (0,0) {transactions et\\hiérarchie};
\node[box,fill=omThm!10] (hi) at (0,-1.4) {historique et\\scénarios de stress};
\node[box] (sc) at (3.4,-1.4) {génération de\\scénarios};
\node[box] (rv) at (3.4,0.7) {réévaluation~:\\complète, grille,\\delta--gamma};
\node[box,fill=omDef!10] (pl) at (6.8,0) {tableau des P\&L\\scénarios $\times$ transactions};
\node[box] (ag) at (10.2,0.7) {agrégation\\par nœud};
\node[box] (rp) at (10.2,-1.0) {VaR, ES, Euler,\\limites, rapports};
\node[box,dashed] (pr) at (3.4,2.5) {bibliothèque de pricing\\(livre~5)};
\draw[->,thick] (md) -- (rv);
\draw[->,thick] (tr) -- (rv);
\draw[->,thick] (hi) -- (sc);
\draw[->,thick] (sc) -- (rv);
\draw[<->,thick,dashed] (pr) -- (rv);
\draw[->,thick] (rv) -- (pl);
\draw[->,thick] (pl) -- (ag);
\draw[->,thick] (ag) -- (rp);
\end{tikzpicture}

\omcaption{Les étapes du \omterm{def:rc:build-a-risk-engine:engine}{moteur de risque}. La réévaluation appelle la bibliothèque de pricing par son interface fixe~;
tout ce qui suit travaille sur le tableau des P\&L par scénario et par transaction. Schéma indicatif.}\label{fig:rc:build-a-risk-engine:arch}
\end{omfigure}

La bibliothèque du livre~5 connaît les données de marché, les instruments, les modèles et les moteurs, et expose \texttt{price},
\texttt{sensitivities}, \texttt{reprice} et \texttt{price\_batch} sur des instantanés immuables décalés par des \omterm{def:rc:market-risk-measures:factor}{facteurs de risque} nommés. Ce livre y ajoute
ses instruments de taux~: une courbe zéro à \omterm{def:rc:curve-construction:pillar}{pilier} qui se décale \omterm{def:rc:curve-construction:pillar}{pilier} par \omterm{def:rc:curve-construction:pillar}{pilier}, des swaps, et des swaptions sous le modèle normal
du chapitre~4, enregistrés avec \texttt{register} de la bibliothèque
afin que ses appels génériques les price comme n'importe quel autre instrument.

\omcode{code/firm/riskengine/firm_riskengine.py}{131}{149}{Le moteur de swaptions et l'enregistrement des instruments de taux auprès de la bibliothèque de pricing.}\label{lst:rc:build-a-risk-engine:register}

\section{Génération des scénarios}

Chaque \omterm{def:rc:stress-testing-and-scenarios:historical}{scénario historique} est un ensemble de chocs sur des \omterm{def:rc:market-risk-measures:factor}{facteurs de risque} nommés~: ici les huit \omterm{def:rc:curve-construction:pillar}{piliers} de la courbe du Trésor
(les taux au pair traités comme des taux zéro, une simplification) et le taux de change euro--dollar au comptant. Le marché de test du moteur
est celui du 23 septembre 2026~: des taux allant de 4.49~\% à un an à 5.40~\% à trente ans,
et EUR/USD à 1.1411. Il construit 250 scénarios d'un jour sur l'année la plus récente et 250 scénarios de dix jours qui se chevauchent~;
les mouvements sur dix jours atteignent 48 points de base sur un \omterm{def:rc:curve-construction:pillar}{pilier} et 2.8~\% sur le taux de change.
Les scénarios filtrés ou soumis à un stress (chapitres~21 et~22) s'insèrent au même endroit.

\section{Réévaluation~: complète, par grille et fondée sur la sensibilité}

\begin{definition}[Réévaluation complète, grille de réévaluation]\label{def:rc:build-a-risk-engine:reval}
La \emph{réévaluation complète}\index{réévaluation complète} réévalue chaque position avec son modèle de pricing dans chaque
scénario. Une \emph{grille de réévaluation}\index{grille de réévaluation} réévalue chaque position sur une petite grille de chocs
de chacun des \omterm{def:rc:market-risk-measures:factor}{facteurs de risque} dont elle dépend, une seule fois, puis calcule le P\&L de chaque scénario par interpolation sur la grille
et en ajoutant les contributions des facteurs.
\end{definition}

La \omterm{def:rc:build-a-risk-engine:reval}{réévaluation complète} coûte un appel de pricing par transaction et par scénario~; le moteur l'exécute par
\texttt{price\_batch} de la bibliothèque, qui signale une transaction en échec au lieu d'interrompre le calcul. La grille coûte un nombre fixe par transaction
et par facteur, quel que soit le nombre de scénarios~; l'\omterm{def:rc:market-risk-measures:dg}{approximation delta--gamma} du chapitre~21 en coûte deux par
transaction et par facteur. La grille capture la non-linéarité de chaque facteur sur toute la plage des mouvements~; comme
l'\omterm{def:rc:market-risk-measures:dg}{approximation delta--gamma} diagonale, elle néglige les effets croisés entre facteurs.

\omcode{code/firm/riskengine/firm_riskengine.py}{206}{222}{Réévaluation par grille~: sept points par facteur, couvrant le plus grand mouvement du scénario.}\label{lst:rc:build-a-risk-engine:grid}

\begin{example}[Coût et erreur sur un livre de 1\,000 transactions]\label{ex:rc:build-a-risk-engine:methods}
Le livre de test compte 600 swaps, 200 swaptions et 200 options EUR/USD, dont 89 arrivent à échéance dans
le mois (71 vendues). La \omterm{def:rc:build-a-risk-engine:reval}{réévaluation complète} sous 250 scénarios exige 251\,000 appels de pricing, la grille 50\,200 et
delta--gamma 17\,400. Avec des scénarios d'un jour, la VaR à 99~\% de la grille reste à moins de 1.59~\% de la \omterm{def:rc:build-a-risk-engine:reval}{réévaluation complète} à chaque nœud, et delta--gamma à moins de 0.74~\%.
Avec des scénarios de dix jours, delta--gamma surestime de 18.45~\% la VaR du desk de change à court terme, et les deux approximations
manquent la VaR des options de taux de plus de 8.6~\%
(\cref{fig:rc:build-a-risk-engine:errors}).
\end{example}

\begin{omfigure}
\begin{tikzpicture}
\begin{axis}[width=12cm,height=5.8cm,ybar,bar width=6pt,
  symbolic x coords={Firm,Rates,FX,Swaps,Rate options,FX short,FX long},xtick=data,xticklabels={Firme,Taux,Change,Swaps,{Options de taux},{Change court},{Change long}},
  x tick label style={font=\footnotesize,rotate=25,anchor=north east},
  ylabel={Erreur de VaR à 10 jours (\%)},
  tick label style={font=\small},label style={font=\small},
  ymin=-5,ymax=20,grid=major,grid style={gray!25},
  legend columns=3,legend style={font=\footnotesize,at={(0.5,1.03)},anchor=south,draw=none,fill=none,/tikz/every even column/.append style={column sep=8pt}},
  table/col sep=comma]
\addplot[fill=omDef!60,draw=omDef] table[x=node,y=grid]{figdata/rates-credit-risk/29-build-a-risk-engine/errors_10d.csv};
\addplot[fill=omThm!60,draw=omThm] table[x=node,y=dg]{figdata/rates-credit-risk/29-build-a-risk-engine/errors_10d.csv};
\addplot[fill=gray!50,draw=gray] table[x=node,y=plan]{figdata/rates-credit-risk/29-build-a-risk-engine/errors_10d.csv};
\legend{grille,delta--gamma,plan de réévaluation}
\end{axis}
\end{tikzpicture}

\omcaption{Erreur relative de la VaR à 10 jours, à 99~\%, par rapport à la \omterm{def:rc:build-a-risk-engine:reval}{réévaluation complète}, par nœud de la \omterm{def:rc:build-a-risk-engine:hierarchy}{hiérarchie de risque}.
La méthode delta--gamma échoue sur les options à courte échéance, dont le gamma varie lors d'un mouvement de dix jours~; les deux approximations
échouent sur les swaptions, dont la valeur dépend du produit des mouvements à différents \omterm{def:rc:curve-construction:pillar}{piliers}~; le plan qui réévalue les swaptions en totalité
reste dans 0.76~\% partout. Données~: le tutoriel du chapitre, sur les rendements du Trésor américain
et les taux de référence de la BCE.}\label{fig:rc:build-a-risk-engine:errors}
\end{omfigure}

Les deux échecs ont des causes différentes (la \cref{fig:rc:build-a-risk-engine:fx}). Les options FX à courte échéance
dépendent d'un seul facteur, et leur P\&L sur un mouvement de 2.8~\% est loin d'une parabole~: la grille, qui
réévalue à la taille du mouvement, la suit~; le développement de Taylor ne le fait pas. Les swaptions dépendent du
taux de swap à terme, combinaison de plusieurs \omterm{def:rc:curve-construction:pillar}{piliers}, et leur convexité tient aux produits croisés des
mouvements de \omterm{def:rc:curve-construction:pillar}{piliers} qu'aucune méthode par facteur ne voit.

\begin{omfigure}
\begin{tikzpicture}
\begin{axis}[width=11cm,height=6cm,
  xlabel={Mouvement EUR/USD sur 10 jours (\%)},ylabel={P\&L du desk (millions d'USD)},
  tick label style={font=\small},label style={font=\small},grid=major,grid style={gray!25},
  legend cell align=left,legend style={font=\footnotesize,at={(0.5,0.04)},anchor=south},
  table/col sep=comma]
\addplot[only marks,mark=o,mark size=1.4pt,draw=gray] table[x=move,y=full]{figdata/rates-credit-risk/29-build-a-risk-engine/fx_short.csv};
\addplot[only marks,mark=x,mark size=1.4pt,draw=omDef] table[x=move,y=grid]{figdata/rates-credit-risk/29-build-a-risk-engine/fx_short.csv};
\addplot[only marks,mark=+,mark size=1.4pt,draw=omThm] table[x=move,y=dg]{figdata/rates-credit-risk/29-build-a-risk-engine/fx_short.csv};
\legend{réévaluation complète,grille,delta--gamma}
\end{axis}
\end{tikzpicture}

\omcaption{P\&L du desk FX à courte échéance (surtout des options vendues) pour chaque scénario de dix jours, en regard du
mouvement EUR/USD du scénario. La grille se superpose à la \omterm{def:rc:build-a-risk-engine:reval}{réévaluation complète}~; la parabole delta--gamma surestime les
pertes des grands mouvements. Données~: le tutoriel du chapitre.}\label{fig:rc:build-a-risk-engine:fx}
\end{omfigure}

\begin{method}[Un plan de réévaluation]\label{met:rc:build-a-risk-engine:plan}
Mesurer l'erreur de chaque approximation par rapport à la \omterm{def:rc:build-a-risk-engine:reval}{réévaluation complète}, nœud par nœud, sur les scénarios qui comptent
(l'horizon le plus long, la période de stress). Attribuer à chaque desk la méthode la moins coûteuse dans le budget d'erreur,
la \omterm{def:rc:build-a-risk-engine:reval}{réévaluation complète} là où aucune ne convient, et relancer la comparaison lorsque le portefeuille ou les scénarios changent.
\end{method}

\omcode{code/firm/riskengine/firm_riskengine.py}{245}{256}{Un plan de réévaluation~: chaque transaction réévaluée par la méthode de son nœud prévu le plus profond.}\label{lst:rc:build-a-risk-engine:plan}

\section{Agrégation, hiérarchies et limites}

\begin{definition}[Hiérarchie de risque, limite de risque]\label{def:rc:build-a-risk-engine:hierarchy}
Une \emph{hiérarchie de risque}\index{hiérarchie de risque} est l'arborescence le long de laquelle positions et risques
sont agrégés~: les transactions dans les books, les books dans les desks, les desks dans les lignes métier et l'entreprise. Une
\emph{limite de risque}\index{limite de risque} est un plafond sur une mesure de risque (VaR, sensibilité, perte de stress) à un
nœud de la hiérarchie, dont le dépassement déclenche une escalade.
\end{definition}

La VaR et l'ES ne sont pas additives~: la mesure de chaque nœud provient de son propre vecteur de P\&L, somme des
colonnes du tableau de P\&L de ses transactions. C'est pourquoi le moteur conserve le tableau, et non des chiffres par desk. Les contributions d'Euler
(chapitre~20) répartissent l'ES d'un nœud entre ses enfants, de sorte que les parties s'additionnent~: la perte moyenne de chaque enfant dans les
scénarios les plus défavorables du parent.

\begin{example}[Risque à dix jours par nœud]\label{ex:rc:build-a-risk-engine:report}
Avec la \omterm{def:rc:build-a-risk-engine:reval}{réévaluation complète}, la VaR de l'entreprise à 10 jours et à 99~\% vaut 53.53~millions de USD, et son ES à 97.5~\%, 50.89~millions de USD.
L'ES propre de l'activité FX est de 14.63~millions de USD, mais sa contribution d'Euler à celle de l'entreprise n'est que de 3.56~millions de USD,
contre 47.34~millions de USD pour les taux~: la queue de distribution de l'entreprise est une queue de taux. La VaR du desk FX à courte échéance, 14.58~millions de USD,
dépasse sa limite de 12~millions de USD, et le moteur le signale.
\end{example}

\omcode{code/firm/riskengine/cpp/firm_riskengine.hpp}{49}{72}{Le noyau C++20~: VaR, ES et contributions d'Euler du vecteur de P\&L d'un nœud (le jumeau Rust est identique ligne à ligne).}\label{lst:rc:build-a-risk-engine:kernel}

\begin{dated}{2026-09}{Agrégation des données de risque}\label{dat:rc:build-a-risk-engine:bcbs239}
Les principes du Comité de Bâle pour l'agrégation efficace des données de risque et la communication des risques (connus
sous le nom de BCBS~239, janvier 2013) fixent quatorze principes, onze pour les banques (gouvernance, architecture des données,
exactitude, exhaustivité, rapidité, adaptabilité, et cinq sur les rapports) et trois pour les superviseurs. Les banques
d'importance systémique mondiale désignées en 2011 ou 2012 devaient se conformer avant janvier 2016. Le dernier
rapport d'avancement du Comité (novembre 2023) évalue 31 de ces banques~: deux seulement étaient pleinement conformes à
l'ensemble des principes, et aucun principe n'était pleinement mis en œuvre par toutes les banques~; sa lettre d'information de janvier 2026 citait
la traçabilité des données, la gouvernance et les données transfrontières comme défis persistants. Le guide de la BCE de mai 2024
rappelait qu'aucun des 25 établissements importants examinés en 2016 ne suivait pleinement les principes, et fixait
sept domaines clés d'attentes de supervision.
\end{dated}

\section{Tester un moteur de risque}

Un \omterm{def:rc:build-a-risk-engine:engine}{moteur de risque} se teste comme une chaîne. La bibliothèque de pricing a ses propres tests d'acceptation. Les
instruments de taux sont testés en formules fermées (un swap au taux au pair vaut zéro, swaption payeuse moins
swaption receveuse est le swap à terme). Les approximations sont testées par rapport à la \omterm{def:rc:build-a-risk-engine:reval}{réévaluation complète} sur de petits mouvements, où
elles doivent concorder. Le noyau est testé sur un jeu de test commun à ses versions Python, C++20 et Rust, et sur
l'identité selon laquelle les contributions d'Euler s'additionnent à l'ES du parent. Compter les appels de pricing transforme le
délai en un nombre que les tests peuvent vérifier.

\section{Tutoriel~: le calcul d'une nuit}\label{tut:rc:build-a-risk-engine:tutorial}

\begin{tutorial}
\textbf{Objectif.} Lancer le moteur sur le portefeuille de 1\,000 transactions sous des \omterm{def:rc:stress-testing-and-scenarios:historical}{scénarios historiques} à un jour et à dix jours~;
produire la VaR, l'ES et les contributions d'Euler par nœud~; comparer la \omterm{def:rc:build-a-risk-engine:reval}{réévaluation complète}, par grille et delta--gamma, en coût et en erreur
~; construire un plan de réévaluation. \textbf{État final~:} les chiffres des
\cref{ex:rc:build-a-risk-engine:methods,ex:rc:build-a-risk-engine:report} et les trois graphiques.
\begin{tutsteps}
\item \textbf{Marché et portefeuille}~: \texttt{market()}, \texttt{book()}.
\item \textbf{Scénarios}~: \texttt{scenarios(1)}, \texttt{scenarios(10)}.
\item \textbf{Réévaluation}~: \texttt{revaluations(h)}, tableaux de P\&L et appels de pricing par méthode.
\item \textbf{Rapports}~: \texttt{reports(h)}, \texttt{errors(h)}, \texttt{euler()}, \texttt{limits()}~;
\texttt{fig\_rc\_riskengine.py}.
\end{tutsteps}
\textbf{Pour aller plus loin.} Ajouter une grille à décalage parallèle pour les swaptions~; utiliser des scénarios de stress issus de
2020 ou 2022~; exécuter le noyau en C++ sur le tableau de P\&L écrit sur disque.
\end{tutorial}

\section{Construction~: le moteur de risque}\label{bld:rc:build-a-risk-engine:riskengine}

\begin{build}
\bfield{Rôle} Le calcul de risque nocturne de l'entreprise~: scénarios, réévaluation via la bibliothèque de pricing,
agrégation le long de la hiérarchie, VaR, ES, contributions d'Euler et limites.
\bfield{Interface} \texttt{PillarCurve}, \texttt{IRSwap}, \texttt{Swaption}, \texttt{NormalRates}
\bfield{Interface} (enregistrés auprès de la bibliothèque de pricing)~; \texttt{Trade}~; \texttt{historical\_scenarios}~;
\texttt{full\_revaluation}, \texttt{grid\_revaluation}, \texttt{delta\_gamma\_revaluation},
\texttt{planned\_revaluation}, \texttt{Counter}~; \texttt{aggregate}, \texttt{risk\_report},
\texttt{euler\_es}, \texttt{check\_limits}~; noyaux C++20 et Rust.
\bfield{Règles} Les prix ne viennent que de la bibliothèque de pricing~; le tableau de P\&L est conservé~; toute approximation est
mesurée par rapport à la \omterm{def:rc:build-a-risk-engine:reval}{réévaluation complète}~; une transaction en échec n'arrête jamais le calcul.
\bfield{Tests d'acceptation} \texttt{code/firm/riskengine/tests/}, \texttt{cpp/}, \texttt{rust/}~: formes fermées des
instruments de taux, chocs sur les \omterm{def:rc:curve-construction:pillar}{piliers}, comptage des appels de pricing de chaque méthode, concordance sur les petits mouvements, plans, limites, et le
jeu de test commun du noyau.
\bfield{Pour aller plus loin} Grilles de \omterm{def:rc:rates-risk:crossgamma}{gamma croisé} ou à décalage parallèle~; calculs incrémentaux pour les nouvelles transactions~; réévaluation parallèle~;
scénarios de stress et scénarios filtrés~; traçabilité de chaque chiffre jusqu'à ses transactions et données.
\end{build}

\begin{omsources}
\item Comité de Bâle sur le contrôle bancaire, \emph{Principles for effective risk data aggregation and risk
reporting}, janvier 2013.
\item Comité de Bâle sur le contrôle bancaire, \emph{Progress in adopting the Principles for effective risk data
aggregation and risk reporting}, novembre 2023.
\item Comité de Bâle sur le contrôle bancaire, lettre d'information n°~36 sur l'agrégation des données de risque, janvier 2026.
\item Banque centrale européenne, \emph{Guide on effective risk data aggregation and risk reporting}, mai 2024.
\item US Department of the Treasury, Daily Treasury Par Yield Curve Rates~; Banque centrale européenne, euro foreign
exchange reference rates (données).
\end{omsources}

\section{Exercices}

\begin{exercise}[$\star$]\label{exo:rc:build-a-risk-engine:1}
Compter les appels de pricing d'une \omterm{def:rc:build-a-risk-engine:reval}{réévaluation complète}, d'une grille (sept points) et d'une réévaluation delta--gamma pour un swap et une
option de change sous 250 scénarios, puis vérifier les totaux du livre.
\end{exercise}

\begin{exercise}[$\star$]\label{exo:rc:build-a-risk-engine:2}
Pourquoi le coût de la grille ne dépend-il pas du nombre de scénarios~?
\end{exercise}

\begin{exercise}[$\star$]\label{exo:rc:build-a-risk-engine:3}
Pourquoi le moteur doit-il conserver le tableau de P\&L plutôt que la VaR de chaque desk~?
\end{exercise}

\begin{exercise}[$\star\star$]\label{exo:rc:build-a-risk-engine:4}
Quelle est l'ampleur de la diversification entre les taux et le change dans l'\omterm{def:rc:market-risk-measures:es}{expected shortfall} à dix jours de la banque~?
\end{exercise}

\begin{exercise}[$\star\star$]\label{exo:rc:build-a-risk-engine:5}
Pourquoi la réévaluation delta--gamma surestime-t-elle les pertes du desk à court terme sur les grands mouvements~?
\end{exercise}

\begin{exercise}[$\star\star$]\label{exo:rc:build-a-risk-engine:6}
Pourquoi aucune méthode par facteur ne capture-t-elle le risque à dix jours des swaptions, et comment y remédier à moindre coût~?
\end{exercise}

\begin{exercise}[$\star\star\star$]\label{exo:rc:build-a-risk-engine:7}
\emph{Programmation.} Compter les appels du plan de réévaluation et les confronter à la composition du plan.
\end{exercise}

\begin{exercise}[$\star\star\star$]\label{exo:rc:build-a-risk-engine:8}
\emph{Trouver la faille.} «~La grille a réussi le test d'erreur l'an dernier, elle est donc conservée.~»
\end{exercise}

\section{Problème~: l'échéance de six heures}

\begin{problem}[{Problème du week-end --- l'échéance de six heures}]\label{pb:rc:build-a-risk-engine:1}
Le véritable livre de la banque compte 50\,000 transactions, comme celui du jeu de test, et le traitement par lots peut utiliser 16 cœurs de deux
heures à six heures du matin, et un appel de pricing coûte en moyenne 20 millisecondes (hypothèses). La mesure de risque est la VaR à 99~\% à
dix jours sous 250 scénarios, et la VaR de chaque nœud doit se situer à moins de 2~\% de la \omterm{def:rc:build-a-risk-engine:reval}{réévaluation complète}.

\medskip\noindent\textbf{Partie I --- Coût.}
\begin{enumerate}
\item Donner les appels de pricing de chaque méthode sur le jeu de test.
\item Les convertir en heures pour le véritable livre.
\item Quelles méthodes tiennent dans la fenêtre~?
\item De quoi dépend le coût de la grille, et celui de la \omterm{def:rc:build-a-risk-engine:reval}{réévaluation complète}~?
\item Combien de scénarios la \omterm{def:rc:build-a-risk-engine:reval}{réévaluation complète} pourrait-elle se permettre dans la fenêtre~?
\end{enumerate}

\medskip\noindent\textbf{Partie II --- Erreur.}
\begin{enumerate}[resume]
\item Donner l'erreur maximale sur les nœuds de chaque approximation à un jour.
\item Et à dix jours.
\item Quel desk fait échouer delta--gamma, et pourquoi~?
\item Quel desk fait échouer la grille, et pourquoi~?
\item Pourquoi l'erreur croît-elle avec l'horizon~?
\end{enumerate}

\medskip\noindent\textbf{Partie III --- Le plan.}
\begin{enumerate}[resume]
\item Construire un plan dans le budget à partir des erreurs sur les nœuds.
\item Donner son coût en appels et en heures.
\item Donner son erreur maximale sur les nœuds.
\item Le plan pourrait-il utiliser delta--gamma pour les options de change à long terme~? Pourquoi est-ce sûr~?
\item Que faut-il surveiller pour que le plan reste valide~?
\end{enumerate}

\medskip\noindent\textbf{Partie IV --- Jugement.}
\begin{enumerate}[resume]
\item Quels dépassements de limites le rapport fait-il apparaître~?
\item Que révèle la décomposition d'Euler sur la queue de distribution de la banque~?
\item Que demanderaient les principes de ponctualité et d'exactitude de BCBS~239 à ce traitement~?
\item Énoncer le \emph{résultat nommé}~: les heures et l'erreur maximale de VaR de la \omterm{def:rc:build-a-risk-engine:reval}{réévaluation complète}, de la grille, de delta--gamma et de la
réévaluation planifiée, ainsi que la méthode la moins coûteuse dans le budget de 2~\%.
\item Pourquoi la méthode la moins coûteuse dans le budget n'est-elle pas un choix figé~?
\end{enumerate}
\end{problem}

\section{Questions d'entretien}

\begin{interviewq}[$\star$ \iqroles{risque}]\label{iq:rc:build-a-risk-engine:1}
Décrivez l'architecture d'un \omterm{def:rc:build-a-risk-engine:engine}{moteur de risque}.
\end{interviewq}

\begin{interviewq}[$\star\star$ \iqroles{développeur}]\label{iq:rc:build-a-risk-engine:2}
Comment feriez-vous tenir un calcul de VaR de nuit sur 50\,000 transactions dans une fenêtre de quatre heures~?
\end{interviewq}

\begin{interviewq}[$\star\star$ \iqroles{risque, chercheur}]\label{iq:rc:build-a-risk-engine:3}
Comparez la \omterm{def:rc:build-a-risk-engine:reval}{réévaluation complète}, les grilles et la réévaluation fondée sur les sensibilités.
\end{interviewq}

\begin{interviewq}[$\star\star$ \iqroles{risque}]\label{iq:rc:build-a-risk-engine:4}
Pourquoi les chiffres de VaR ne sont-ils pas additifs d'un desk à l'autre, et comment les répartissez-vous~?
\end{interviewq}

\begin{interviewq}[$\star\star\star$ \iqroles{développeur, risque}]\label{iq:rc:build-a-risk-engine:5}
Comment testez-vous un \omterm{def:rc:build-a-risk-engine:engine}{moteur de risque}~?
\end{interviewq}

\begin{interviewq}[$\star\star\star$ \iqroles{banque, risque}]\label{iq:rc:build-a-risk-engine:6}
Que demande BCBS~239, et pourquoi les banques peinent-elles à s'y conformer~?
\end{interviewq}
